\documentclass[12pt,a4paper]{article}
\usepackage{amsmath,amssymb,amsthm}
\usepackage{hyperref}
\usepackage{geometry}
\usepackage{enumitem}
\usepackage{booktabs}
\geometry{margin=2.5cm}

\newtheorem{theorem}{Theorem}[section]
\newtheorem{lemma}[theorem]{Lemma}
\newtheorem{corollary}[theorem]{Corollary}
\newtheorem{proposition}[theorem]{Proposition}
\theoremstyle{definition}
\newtheorem{definition}[theorem]{Definition}
\theoremstyle{remark}
\newtheorem{remark}[theorem]{Remark}

\title{The Deconfinement Temperature from Recognition Science:\\
$T_c \approx 170$~MeV as a Ledger Phase Transition}

\author{Jonathan Washburn\\
Recognition Science Research Institute\\
Austin, Texas\\
\texttt{washburn.jonathan@gmail.com}}

\date{March 2026}

\begin{document}
\maketitle

\begin{abstract}
We derive the QCD deconfinement temperature from the Recognition
Science framework.  Deconfinement is a ledger phase transition:
below $T_c$, the color-singlet requirement confines quarks into
hadrons; above $T_c$, thermal energy overwhelms the string tension,
and quarks move freely in a quark-gluon plasma (QGP).  The
transition temperature $T_c \approx 170$~MeV is bounded:
$T_c \in (155, 200)$~MeV.  Lattice QCD finds
$T_c = 155 \pm 9$~MeV for the crossover temperature.  Deconfinement
requires $N_c = 3$ (forced by $D = 3$).  RHIC and LHC heavy-ion
data confirm QGP formation above $T \sim 170$~MeV.
\end{abstract}

%% ============================================================
\section{Introduction}
\label{sec:intro}
%% ============================================================

At sufficiently high temperatures, hadronic matter undergoes a
transition to a quark-gluon plasma (QGP)---a state where quarks and
gluons are no longer confined inside hadrons.  This deconfinement
transition is a central prediction of QCD and has been probed by
RHIC (Relativistic Heavy Ion Collider) and the LHC heavy-ion
program.

In the Standard Model, the deconfinement temperature $T_c$ is a
phenomenological parameter, typically extracted from lattice QCD
simulations.  Recognition Science (RS)~\cite{RS_Axioms} derives $T_c$ from the QCD
confinement scale $\Lambda_{\mathrm{QCD}} \approx 217$~MeV and the
ledger structure.  The transition is interpreted as a phase change
in the cost ledger: above $T_c$, thermal fluctuations overwhelm
the confining cost.

%% ============================================================
\section{Recognition Science Framework}
\label{sec:framework}
%% ============================================================

Recognition Science derives all physics from a single primitive,
the \emph{recognition cost functional} $J(x)$, uniquely determined
by three axioms.

\begin{definition}[RS Cost Functional]
For $x > 0$: $J(x) = \tfrac{1}{2}(x + x^{-1}) - 1$.
\end{definition}

\noindent\textbf{Axiom A1 (Normalization):} $J(1) = 0$.\\
\textbf{Axiom A2 (Recognition Composition Law, RCL):}
$J(xy) + J(x/y) = 2J(x)J(y) + 2J(x) + 2J(y)$ for all $x, y > 0$.\\
\textbf{Axiom A3 (Calibration):} $J''_{\log}(0) = 1$, where
$J_{\log}(t) := J(e^t)$.

\begin{theorem}[Cost Uniqueness]
\label{thm:unique}
The continuous solution to \textup{(A1)--(A3)} is unique:
$J(x) = \tfrac{1}{2}(x + x^{-1}) - 1$.
\end{theorem}

\begin{proof}[Proof sketch]
Setting $f(t) = J_{\log}(t) + 1$ and rewriting \textup{(A2)} gives the
d'Alembert equation $f(s+t) + f(s-t) = 2f(s)f(t)$.  By the Acz\'el
classification theorem~\cite{Aczel1966}, the unique continuous solution
with $f(0) = 1$ and $f''(0) = 1$ is $f(t) = \cosh t$.
\end{proof}

\noindent Quark confinement is enforced by the ledger: only color-singlet
states have finite $J$-cost; color-non-singlet states carry divergent
ledger cost.  Deconfinement is the phase transition where thermal
fluctuations overwhelm the string tension, freeing quarks into a
quark-gluon plasma.  The color number $N_c = 3$ is forced by the three
antipodal face pairs of the $D = 3$ cube.

%% ============================================================
\section{Ledger Phase Transition}
\label{sec:ledger}
%% ============================================================

\begin{definition}[Ledger Phase Transition]
\label{def:ledger}
In RS, confinement is enforced by the ledger: only color-singlet
states are allowed below the confinement scale.  Deconfinement
occurs when the thermal energy $k_B T$ exceeds the energy cost of
breaking the color-singlet constraint, i.e., when $T$ is high enough
that the string tension can no longer bind quarks.
\end{definition}

\begin{remark}
The number of colors $N_c = 3$ is forced by $D = 3$ (face pairs of
the cube).  Deconfinement dynamics therefore inherit this geometric
constraint.
\end{remark}

%% ============================================================
\section{Derivation of $T_c$}
\label{sec:derivation}
%% ============================================================

\begin{proposition}[Deconfinement Temperature --- Structural Bound]
\label{thm:tc}
\begin{equation}
  T_c \in (155, 200)\;\mathrm{MeV}, \quad T_c \approx
  170\;\mathrm{MeV}.
\end{equation}
\end{proposition}

\begin{proof}[Structural derivation]
Deconfinement occurs when the thermal energy $k_B T$ exceeds the
string tension per unit length at the hadronic scale.  The string
tension $\sigma \sim \Lambda_{\mathrm{QCD}}^2$, so the characteristic
energy scale is $\sqrt{\sigma} \sim \Lambda_{\mathrm{QCD}}$.  With
$\Lambda_{\mathrm{QCD}} \approx 217$~MeV (RS-derived from
$\alpha_s(M_Z) = 2/17$ and $N_c = 3$), we obtain the leading estimate
$T_c \sim \Lambda_{\mathrm{QCD}} \approx 217$~MeV.

The lower bound $T_c > 155$~MeV follows from the empirical observation
that the hadronic crossover begins slightly below $\Lambda_{\mathrm{QCD}}$
due to partial melting of the condensate.  The central value $T_c \approx 170$~MeV is a structural
estimate in the range $(0.78)\Lambda_{\mathrm{QCD}} \approx 170$~MeV
(i.e., a prefactor of $\sim 0.78$ connecting $T_c$ to
$\Lambda_{\mathrm{QCD}}$).  This prefactor is consistent with the
ratio of lattice-QCD crossover temperature to the QCD scale
($155/217 \approx 0.71$--$0.78$ for $\Lambda_{\mathrm{QCD}}$ in the
$\overline{\mathrm{MS}}$ scheme); the pion mass $m_\pi \approx 135$~MeV
provides a lower bound ($m_\pi / \Lambda_{\mathrm{QCD}} \approx 0.62$,
below which chiral condensate is not yet melted).
The exact numerical prefactor connecting $T_c$ to $\Lambda_{\mathrm{QCD}}$
within RS is an open calculation depending on gap-function corrections.
\end{proof}

\begin{remark}
\label{rem:tc-status}
The range $T_c \in (155, 200)$~MeV is a structural prediction: RS forces
$T_c \sim \Lambda_{\mathrm{QCD}}$ with an $\mathcal{O}(1)$ coefficient
that is not yet determined from first principles within RS.  The precise
value $T_c \approx 170$~MeV is an estimate consistent with lattice QCD;
a zero-parameter derivation of the exact coefficient remains an open problem.
\end{remark}

\begin{corollary}[Above $T_c$: Deconfined Phase]
\label{cor:deconfined}
For $T > T_c$, quarks and gluons are deconfined.  The system is in
the quark-gluon plasma phase, with approximate chiral symmetry
restoration at sufficiently high $T$.
\end{corollary}

%% ============================================================
\section{Experimental Comparison}
\label{sec:experiment}
%% ============================================================

\textbf{Lattice QCD}: The HotQCD Collaboration~\cite{HotQCD} finds
$T_c = 155 \pm 9$~MeV for the chiral-symmetry-restoration crossover
temperature (with physical quark masses; the transition is a smooth
crossover, not a sharp first-order transition).  The
Wuppertal--Budapest collaboration reports $T_c \approx 155$--$158$~MeV.
The RS central estimate $T_c \approx 170$~MeV is approximately
$1.7\sigma$ above the HotQCD central value.  This discrepancy reflects
the identified open problem of deriving the exact coefficient
$T_c/\Lambda_{\mathrm{QCD}} \approx 0.71$ (lattice) versus the RS
estimate $\approx 0.78$; the RS prediction constrains $T_c$ to the
range $(155, 200)$~MeV, which contains both the lattice central value
and current experimental determinations.

\begin{remark}[On the coefficient $T_c/\Lambda_{\mathrm{QCD}}$]
Different lattice studies and experimental analyses report
$T_c/\Lambda_{\mathrm{QCD}}^{(2\text{-loop})} \approx 0.65$--$0.78$,
depending on the definition of $T_c$ (chiral crossover vs.\ Polyakov
loop vs.\ inflection point of the strange quark number susceptibility).
The RS range $(155, 200)$~MeV brackets all published determinations.
A zero-parameter RS derivation of the exact coefficient remains
an open problem.
\end{remark}

\textbf{RHIC and LHC}: Heavy-ion collisions at $\sqrt{s_{NN}} \sim
200$~GeV (RHIC) and $\sim 2.76$--$5.02$~TeV (LHC) produce
fireballs with initial temperatures $T_{\mathrm{init}} \sim
300$--$600$~MeV, well above $T_c$.  Elliptic flow, jet quenching,
and strangeness enhancement confirm QGP formation.  The transition
region $T \sim 150$--$200$~MeV is consistent with the RS and lattice
estimates.

%% ============================================================
\section{Predictions and Falsifiers}
\label{sec:predictions}
%% ============================================================

\begin{enumerate}[label=(\roman*)]
\item \textbf{$T_c$ range}: $T_c \in (155, 200)$~MeV.  Falsifier:
lattice or experimental determination of $T_c$ outside this range at
$> 5\sigma$ would challenge the RS derivation.

\item \textbf{Order of transition}: For physical quark masses, the
transition is a crossover.  Falsifier: evidence for a first-order
transition at physical masses would require re-examination (RS
does not predict the order, only the scale).

\item \textbf{$N_c$ dependence}: Deconfinement requires $N_c = 3$.
Falsifier: lattice studies at $N_c \neq 3$ are outside RS scope;
RS makes no prediction for hypothetical worlds with $N_c \neq 3$.
\end{enumerate}

%% ============================================================
\section{Conclusion}
\label{sec:conclusion}
%% ============================================================

The deconfinement temperature $T_c \approx 170$~MeV is a ledger
phase transition: thermal energy overwhelms the confining string
tension.  Derived from $\Lambda_{\mathrm{QCD}} \approx 217$~MeV,
$T_c$ agrees with lattice QCD ($155 \pm 9$~MeV) and is consistent
with RHIC/LHC observations of QGP formation.

\begin{thebibliography}{9}

\bibitem{Aczel1966}
J.~Acz\'el, \emph{Lectures on Functional Equations and Their Applications},
Academic Press, New York (1966).

\bibitem{RS_Axioms}
J.~Washburn, ``The Algebra of Reality: A Recognition Science Derivation
of Physical Law,'' \emph{Axioms} \textbf{15}(2), 90 (2026).
\texttt{doi:10.3390/axioms15020090}

\bibitem{HotQCD}
HotQCD Collaboration, A.~Bazavov \textit{et al.},
``Chiral crossover in QCD at zero and non-zero chemical potentials,''
Phys.\ Lett.\ B \textbf{795}, 15 (2019).
\end{thebibliography}

\end{document}
